% Generated by paper/scripts/generate_tables.py; do not edit by hand.
\begin{table}[!htbp]
\centering
\scriptsize
\begin{tabular}{@{}L{0.24\linewidth}L{0.68\linewidth}@{}}
\toprule
Calibration area & Use in Adaptive Bridge v1 \\
\midrule
Lobby venue shifting & Anti-capture rows with audit, machine-readable logs, venue-shift detection, or sanctions reduce lobby adaptation and capture pressure; narrow caps and budget tools keep a substitution-risk caveat. \\
Emergency powers & Emergency-integrity rows reduce rights-threat inflow, court adaptation, court-curbing pressure, and review backlog only when review is fast, reason-giving, and anti-evasion oriented. \\
Sequencing and transition costs & Sequencing readiness reduces transition load, agency gaps, and delivery bottlenecks; complex one-shot redesigns carry higher transition and direct-evidence penalties. \\
Federalism and agency capacity & The bridge now tracks federalism resistance, delivery bottleneck pressure, and recovery capacity separately before combining them into agency implementation gaps. \\
Public feedback & Feedback correction capacity can turn backlash into alignment and policy correction, while weak voice causes faster legitimacy erosion. \\
Citizen agenda setting & Citizen assemblies, random public review panels, and comment-authenticity rules improve agenda correction; initiatives and petition channels also add rights-risk or overload pressure. \\
Comparative evidence strength & Low-direct-evidence families can rank highly synthetically, but the recommendation gate labels them as calibration gray-zone or do-not-recommend rather than provisional recommendations. \\
\bottomrule
\end{tabular}
\caption{Evidence-informed priors added to Adaptive Bridge v1. The table records directional model use only; it is not an empirical fit.}
    \end{table}
